\mode*
\section{Misconceptions we have found}

While tutoring students on labs and assignments, we have repeatedly met
misconceptions that the literature does not document, or documents only in
passing. We treat them as we treated the misconceptions drawn from the
literature: as phenomenographic observations to be analysed through variation
theory. Because we observed them directly rather than through a published study,
they should be read as hypotheses for teaching designs rather than as
established findings.

\subsection{The boundary between a function and the rest of the program}

Several of our observations share a single object of learning: \emph{how a
function communicates with the rest of the program}. The critical aspects are
that a function is invoked explicitly; that values enter it only through the
arguments actually passed; that a value leaves it only through
\mintinline{python}{return}, as a value rather than as a name; and that the
names used inside it are local. Students who have not discerned this boundary
produce a family of related errors, which a contrast on each aspect can separate.

\subsubsection{A function must be invoked}

Some students believe that a defined function will be called \enquote{by itself},
in the right order, without ever being invoked. The critical aspect --- that
definition and invocation are distinct --- becomes visible if we hold the
definition invariant and vary only whether the call is present.

\hfill
\begin{minipage}[t]{0.45\columnwidth}
  \begin{minted}{python}
    def greet():
        print("Hi!")
    # nothing is printed
  \end{minted}
\end{minipage}
\hfill
\begin{minipage}[t]{0.45\columnwidth}
  \begin{minted}[highlightlines={3}]{python}
    def greet():
        print("Hi!")
    greet()
  \end{minted}
\end{minipage}

\subsubsection{Arguments must be passed}

A related belief is that, if a parameter shares its name with a global variable,
the argument need not be passed at the call. Holding the function invariant and
varying only whether the arguments are supplied makes the rule discernible: the
left call raises \mintinline{python}{TypeError}, because the parameters are bound
from the arguments at the call, not from globals that happen to share their
names.

\hfill
\begin{minipage}[t]{0.45\columnwidth}
  \begin{minted}[highlightlines={5}]{python}
    def foo(x, y):
        print(x + y)
    x = 3
    y = 4
    foo()
  \end{minted}
\end{minipage}
\hfill
\begin{minipage}[t]{0.45\columnwidth}
  \begin{minted}[highlightlines={5}]{python}
    def foo(x, y):
        print(x + y)
    x = 3
    y = 4
    foo(x, y)
  \end{minted}
\end{minipage}

\subsubsection{A value leaves only through \texttt{return}, as a value}

The most tangled observations concern \mintinline{python}{return}. Some students
think that, to return a variable's value, that variable must have been an input
argument; others think that \mintinline{python}{return} \enquote{sends} the
variable's \emph{name} to the call site, so that the name can be used there
afterwards. Both miss that \mintinline{python}{return} hands back a \emph{value},
which the caller must capture in a name of its own. We contrast relying on the
name appearing at the call site with capturing the returned value.

\hfill
\begin{minipage}[t]{0.45\columnwidth}
  \begin{minted}[highlightlines={4}]{python}
    def foo():
        x = 3
        return x
    foo()
    print(x)   # NameError
  \end{minted}
\end{minipage}
\hfill
\begin{minipage}[t]{0.45\columnwidth}
  \begin{minted}[highlightlines={4}]{python}
    def foo():
        x = 3
        return x
    x = foo()
    print(x)   # 3
  \end{minted}
\end{minipage}

\subsubsection{A mutable argument is shared, not copied}

A subtler boundary misconception runs the other way: that a list passed to a
function is copied, so that changing it inside the function leaves the caller's
list untouched. In fact the function receives a reference to the \emph{same}
list, so a mutation is visible outside. Holding the call invariant and varying
only whether the function \emph{mutates} the list or \emph{rebinds} the name
separates the two: mutation reaches the caller, rebinding does not. This is the
same reference behaviour behind the mutable-default-argument example of
\cref{debugging}.

\hfill
\begin{minipage}[t]{0.45\columnwidth}
  \begin{minted}[highlightlines={2}]{python}
    def f(xs):
        xs.append(4)
    a = [1, 2, 3]
    f(a)
    print(a)   # [1, 2, 3, 4]
  \end{minted}
\end{minipage}
\hfill
\begin{minipage}[t]{0.45\columnwidth}
  \begin{minted}[highlightlines={2}]{python}
    def f(xs):
        xs = [4]
    a = [1, 2, 3]
    f(a)
    print(a)   # [1, 2, 3]
  \end{minted}
\end{minipage}

\subsection{Observations related to the catalogue}

Three further observations restate, from our own teaching, misconceptions
already analysed above, and are best taught alongside them.

\begin{itemize}
  \item That the built-in \mintinline{python}{input()} converts what the user
    types, returning an integer when a number is entered. This is the same
    type misconception we discuss for data types --- \mintinline{python}{input()}
    always returns a string --- and the same contrast (\mintinline{python}{int}
    applied or not) applies.
  \item That a variable's name constrains the type of value it may hold (for
    example, that assigning an integer to a variable named
    \mintinline{python}{my_string} is an error). This is the
    name-has-power-over-value misconception from the functions-and-variables
    chapter.
  \item That a local variable is reachable outside the function in which it is
    defined --- the same scope boundary treated above.
\end{itemize}

A fourth observation concerns repetition: to run a function several times while a
condition holds, some students have the function call \emph{itself} rather than
placing the call in a loop. The contrast between the two is instructive, and
belongs with the treatment of loops and recursion.

\hfill
\begin{minipage}[t]{0.45\columnwidth}
  \begin{minted}{python}
    def foo():
        try:
            x = int(input("Integer:"))
            return x
        except ValueError:
            foo()
  \end{minted}
\end{minipage}
\hfill
\begin{minipage}[t]{0.45\columnwidth}
  \begin{minted}{python}
    def foo():
        while True:
            try:
                return int(input("Integer:"))
            except ValueError:
                continue
  \end{minted}
\end{minipage}

\subsection{Difficulties that are not misconceptions}

Not every difficulty is a misconception. For instance, the word \emph{iterate}
is not a word all students have met before, and needs to be explained in plain
terms; this is a matter of vocabulary rather than a mistaken conception of how
the code behaves.
